\input{sol_pre.tex}
\input{q_2025-12-27.tex}
\begin{document}
\input{sol_head.tex}

\Savol{1}
$88+95+96+a=4\cdot94=376$ $\Rightarrow$ $a=376-279=97$.
\Javob{C}

\Savol{2}
1) $\pi+\sqrt{10}$ — irratsional. \ 2) $\frac{22}7$ — ratsional.\\
3) $t=\sqrt[3]{2-\sqrt5}+\sqrt[3]{2+\sqrt5}$: $t^3=4+3\sqrt[3]{4-5}\cdot t=4-3t$ $\Rightarrow$ $t^3+3t-4=0$ $\Rightarrow$ $t=1$ — ratsional.\\
4) Davriy kasr — ratsional. \ 5) $(\sqrt3-1)+(\sqrt3+1)=2\sqrt3$ — irratsional.\\
Irratsionallar: $2$ ta.
\Javob{A}

\Savol{3}
Şerzod $x$ kun, Anvar $x+5$, Jasur $x-4$ kun: \ $\dfrac1x+\dfrac1{x+5}+\dfrac1{x-4}=\dfrac13$.\\
$x=10$: $\frac1{10}+\frac1{15}+\frac16=\frac{3+2+5}{30}=\frac13$ — mos.
\Javob{A}

\Savol{4}
$400\,000\cdot0{,}7\cdot1{,}1=308\,000$.
\Javob{C}

\Savol{5}
$1^3+2^3+\dots+9^3=(1+2+\dots+9)^2=45^2=2025$ $\Rightarrow$ $\frac{2025}{2025}=1$.
\Javob{A}

\Savol{6}
$a=11$, $b=100$, $a+b=111$: \ $a^4+b^4+(a+b)^4=2\left(a^2+ab+b^2\right)^2$.
\[\sqrt{\frac{2\left(a^2+ab+b^2\right)^2}2}=121+1100+10000=11221.\]
\Javob{B}

\Savol{7}
$\left(1-x^2+2x^3\right)^5=1-5x^2+\dots$ ($x$ li had yöq), \ $(x-2)^3=-8+12x-6x^2+x^3$.\\
$x^2$ oldidagi koeffitsiyent: $1\cdot(-6)+(-5)\cdot(-8)=34$.
\Javob{C}

\Savol{8}
$a_2+a_4=2a_3=16$ $\Rightarrow$ $a_3=8$. \ $a_1a_5=(8-2d)(8+2d)=64-4d^2=28$ $\Rightarrow$ $d=\pm3$.\\
$a_2=8-d=5$ yoki $11$.
\Javob{B}

\Savol{9}
$b_1(1+q+q^2)=91$, $b_1q(q-1)=42$ $\Rightarrow$ $6(1+q+q^2)=13(q^2-q)$ $\Rightarrow$ $7q^2-19q-6=0$ $\Rightarrow$ $q=3$ (ösuvçi).\\
$b_1=\frac{91}{13}=7$, $b_2=21$: \ $S_2=28$.
\Javob{C}

\Savol{10}
$x^2+\frac1{x^2}=36-2=34$. \ $\dfrac{x^4+1}{4x^2}=\dfrac14\left(x^2+\dfrac1{x^2}\right)=\dfrac{34}4=8\dfrac12$.
\Javob{B}

\Savol{11}
$2\left(7+3\sqrt5\right)=14+6\sqrt5=\left(3+\sqrt5\right)^2$ $\Rightarrow$ $\sqrt[6]{2x(7+3\sqrt5)}=\sqrt[3]{\sqrt x\left(3+\sqrt5\right)}$.
\[\sqrt[3]{\sqrt x\left(3+\sqrt5\right)\cdot\sqrt x\left(\sqrt5-3\right)}=\sqrt[3]{x(5-9)}=-\sqrt[3]{4x}.\]
\Javob{C}

\Savol{12}
$\sin40^\circ=2\sin20^\circ\cos20^\circ$:
\[\frac{2\sin20^\circ(1+\cos20^\circ)}{2\sin20^\circ(1-\cos20^\circ)}=\frac{2\cos^210^\circ}{2\sin^210^\circ}=\ctg^210^\circ.\]
\Javob{C}

\Savol{13}
$\sin x\sin y=\frac12\bigl(\cos(x-y)-\cos(x+y)\bigr)=\frac12\left(\cos(x-y)-\frac12\right)=\frac14$ $\Rightarrow$ $\cos(x-y)=1$.\\
$x-y=2\pi k$: \ $x=\frac\pi6+\pi k$, $y=\frac\pi6-\pi k$. Ikkalasi ham $[-2\pi;2\pi]$ da: $k=-1;0;1$ — $3$ ta.
\Javob{C}

\Savol{14}
$x\ge0$ da $x\cdot2^x\ge x$, demak çap tomon $>x$.\\
$x<0$ da $x\cdot2^x>-1$, çap tomon $>8-1=7>x$. Ildiz yöq.
\Javob{C}

\Savol{15}
$t=\log_2x$: \ $t^2+(x-1)t+2x-6=0$, \ $D=(x-1)^2-4(2x-6)=(x-5)^2$.\\
$t=-2$ $\Rightarrow$ $x=\frac14$; \ $t=3-x$: $\log_2x=3-x$ $\Rightarrow$ $x=2$ (yagona, çap ösuvçi, öng kamayuvçi).\\
Yiğindi: $\frac14+2=2{,}25$.
\Javob{C}

\Savol{16}
Ikkinçi tenglamani $3$ ga köpaytirib birinçisini ayiramiz: $7y=35$, $y=5$, $x=7$. \ $x+y=12$.
\Javob{B}

\Savol{17}
$f(x)=x^3-x^2-2x-2002$, \ $f'(x)=3x^2-2x-2=0$ $\Rightarrow$ $x=\frac{1\pm\sqrt7}3$.\\
Lokal maksimum ($x\approx-0{,}55$) da $f<0$ $\Rightarrow$ grafik $Ox$ ni bir marta kesadi: $1$ ta yeçim.
\Javob{B}

\Savol{18}
$x>4$: $x^2-8x+15\le0$ $\Rightarrow$ $3\le x\le5$ $\Rightarrow$ $x\in(4;5]$.\\
$x<4$: $x^2-8x+21\le0$ — $D<0$, yeçim yöq. \ Butun yeçim: $x=5$ — $1$ ta.
\Javob{A}

\Savol{19}
$t=\sqrt{8+x}\ge0$, $x=t^2-8$: \ $x-8t+24=(t-4)^2$, \ $16+x=t^2+8$:
\[|t-4|<\frac{t^2+8}{2(t-5)}.\]
$t<5$ da öng tomon manfiy — yeçim yöq. $t>5$ da: $2(t-4)(t-5)<t^2+8$ $\Rightarrow$ $t^2-18t+32<0$ $\Rightarrow$ $2<t<16$.\\
$5<t<16$ $\Rightarrow$ $25<8+x<256$ $\Rightarrow$ $17<x<248$: \ $248-18=230$ ta butun son.
\Javob{C}

\Savol{20}
Viyet: $x_1\cdot12x_1=3$ $\Rightarrow$ $x_1=\pm\frac12$; \ $13x_1=\frac{10-a}2$ $\Rightarrow$ $a=10-26x_1=-3$ yoki $23$. Yiğindi $20$.
\Javob{A}

\Savol{21}
$u=4-x^3$, $du=-3x^2dx$:
\[\int_0^1\frac{6x^2}{(4-x^3)^2}dx=\int_3^4\frac{2\,du}{u^2}=\left.-\frac2u\right|_3^4=-\frac12+\frac23=\frac16.\]
\Javob{C}

\Savol{22}
Perpendikulyarlik: $a\cdot\left(-\frac12\right)=-1$ $\Rightarrow$ $a=2$. \ $y'=\dfrac{7x/2}{\sqrt{\frac72x^2-7}}=2$ $\Rightarrow$ $\frac{49}4x^2=14x^2-28$ $\Rightarrow$ $x^2=16$, $x=4$ ($y'>0$).\\
$y(4)=\sqrt{49}=7$: urinma $y=2(x-4)+7=2x-1$. \ $a+b=2-1=1$.
\Javob{A}

\Savol{23}
\fig{\begin{tikzpicture}[scale=0.55,font=\small]
\draw (0,0) circle (2.83); \coordinate (A) at (-2,-2); \coordinate (B) at (-2,2); \coordinate (C) at (2,-2);
\draw (B)--(A)--(C); \draw[ax] (B)--(C); \draw[axs] (0,0)--(-2,0); \fill (0,0) circle (2pt);
\node[below left] at (A) {$A$}; \node[above left] at (B) {$B$}; \node[below right] at (C) {$C$}; \node[above right] at (0,0) {$O$};
\node[aux,above] at (-1,0) {$8$};
\end{tikzpicture}}
$AB=AC$, $\angle BAC=90^\circ$ $\Rightarrow$ $BC$ — diametr, $O$ — uning örtasi.\\
$O$ dan $AB$ gaça perpendikulyar — $\triangle ABC$ örta çiziği: $8=\frac{AC}2$ $\Rightarrow$ $AC=16$.
\Javob{A}

\Savol{24}
$\dfrac3{\sqrt{1+x}}=2$ $\Rightarrow$ $1+x=\frac94$, $x=1\frac14$.
\Javob{C}

\Savol{25}
Har bir uçburçak töğri burçakli, markazdagi burçagi $30^\circ$, keyingi uçburçakning gipotenuzasi oldingisining markazga yopişgan kateti: har qadamda uzunlik $\cos30^\circ$ ga köpayadi. $6$ ta uçburçak:
\[x=64\cos^630^\circ=64\cdot\frac{27}{64}=27.\]
\Javob{B}

\Savol{26}
Urinmalar xossasi: \ $AE=\dfrac{AB+AD-BD}2$, \ $AF=\dfrac{AC+AD-DC}2$.
\[EF=AF-AE=\frac{AC-DC-AB+BD}2=\frac{7-3-5+5}2=2.\]
\Javob{B}

\Savol{27}
$\dfrac{AE}{AD}=t$. \ $\triangle ABD$ da $EO\parallel AB$: $EO=4(1-t)$; \ $\triangle BCD$ da $OF\parallel DC$: $OF=12t$.\\
$D$ dan $EF$ gaça masofa $(1-t)h$, $B$ dan $EF$ gaça $th$:
\[\frac12\cdot4(1-t)\cdot(1-t)h=\frac12\cdot12t\cdot th\ \Rightarrow\ 1-t=\sqrt3\,t\ \Rightarrow\ t=\frac1{1+\sqrt3}.\]
$EF=4(1-t)+12t=4+8t=4+4(\sqrt3-1)=4\sqrt3$.
\Javob{A}

\Savol{28}
Muntazam beşburçak burçagi $108^\circ$, $\angle OEF=180^\circ-108^\circ=72^\circ$. \ $\triangle OEF$: $\angle EOF=180^\circ-72^\circ-45^\circ=63^\circ$.\\
$\triangle ODG$: $\angle ODG=54^\circ$, $\angle DOG=63^\circ$ $\Rightarrow$ $\angle DGO=63^\circ$ $\Rightarrow$ $DG=DO=8$.\\
$DG$ va $EH$ — uçdan qarşi tomon örtasigaça kesmalar (simmetriya öqlari), ular teng: $EH=8$.
\Javob{A}

\Savol{29}
Oğirlik markazi: $M\left(\frac{-5+14+3}3;\,\frac{4+5-3}3\right)=(4;\,2)$. \ $OM=\sqrt{16+4}=2\sqrt5$.
\Javob{C}

\Savol{30}
Tekisliklar orasidagi masofa: $h=\sqrt{37^2-12^2}=35$. \ $BD$ proyeksiyasi $\sqrt{125^2-35^2}=\sqrt{14400}=120$.
\Javob{A}

\Savol{31}
$A\cup B=\{a;b;c;d;m;n;k\}$ — $7$ ta element: $2^7=128$.
\Javob{C}

\Savol{32}
$n=5^3=125$; \ $m=4\cdot5^3=500$ (birinçi raqam $0$ bölmaydi). \ $\frac nm=\frac14$.
\Javob{D}

\newpage
\Savol{33}
\fig{\begin{tikzpicture}[scale=0.85,font=\small]
\coordinate (A) at (0,0); \coordinate (B) at (3.2,0); \coordinate (C) at (3.2,3.2); \coordinate (D) at (0,3.2); \coordinate (O) at (1.326,1.874);
\draw (A)--(B)--(C)--(D)--cycle; \draw[ax] (B)--(D); \draw (O) circle (1.874);
\draw[axs] (O)--(1.326,0); \draw[axs] (O)--(3.2,1.874); \fill (O) circle (1.5pt);
\node[below left] at (A) {$A$}; \node[below right] at (B) {$B$}; \node[above right] at (C) {$C$}; \node[above left] at (D) {$D$}; \node[left] at (O) {$O'$};
\node[aux,left] at (1.326,0.8) {$r$};
\end{tikzpicture}}
Şar $B$ uçidagi uçta yoqqa urinadi: $O$ şu yoqlardan $r$ masofada, $O$ ning asosdagi proyeksiyasi $O'$ — $BD$ diagonalda, $BO'=r\sqrt2$, $O$ ning asosdan balandligi $r$.\\
$DD_1\perp(ABCD)$, şuning uçun $O$ dan $DD_1$ gaça masofa $O'D$ ga teng:
\[O'D=BD-r\sqrt2=2-r\sqrt2=r\ \Rightarrow\ r=\frac2{1+\sqrt2}=2\sqrt2-2.\]

\textbf{33.} $r=2\sqrt2-2$ — \textbf{B}.

\textbf{34.} $AA_1D_1D$ yoq $BB_1C_1C$ ga qarama-qarşi: $\sqrt2-r=\sqrt2-2\sqrt2+2=2-\sqrt2$ — \textbf{D}.

\textbf{35.} $OD=\sqrt{O'D^2+r^2}=r\sqrt2=4-2\sqrt2$ — \textbf{F}.
\Javob{33 — B, \ 34 — D, \ 35 — F}

\newpage
\Savol{36}
\fig{\begin{tikzpicture}[scale=0.5,font=\small]
\draw[->] (-3.5,0)--(3.5,0) node[right] {$x$}; \draw[->] (0,-1)--(0,6) node[left] {$y$};
\draw[thick,domain=-2.3:2.3,samples=80] plot (\x,{\x*\x});
\draw[aux] (0,-2) circle (2); \draw[aux,dashed] (0,4.25) circle (2);
\fill[aux] (0,-2) circle (2pt); \fill[aux] (0,4.25) circle (2pt);
\node[aux,right] at (0,-1.6) {\scriptsize$a=2$}; \node[aux,right] at (0.1,4.6) {\scriptsize$a=-\frac{17}4$};
\end{tikzpicture}}
Birinçi tenglama: $x^2+(y+a)^2=4$ — markazi $(0;-a)$, radiusi $2$ aylana. $y=x^2\ge0$ ni qöyamiz:
\[y^2+(2a+1)y+a^2-4=0.\qquad(*)\]
Har bir musbat ildiz $y$ ikkita yeçim ($x=\pm\sqrt y$), $y=0$ esa bitta yeçim beradi.

a) Bitta yeçim: $y=0$ ildiz va musbat ildiz yöq. $a^2=4$: \ $a=2$ da $y^2+5y=0$ — $y=0;-5$, mos; \ $a=-2$ da $y^2-3y=0$ — $y=0;3$, uçta yeçim. Demak $a=2$, yiğindi $2$.

b) Ikkita yeçim: $(*)$ ning bitta musbat ildizi bölişi kerak.\\
— Ildizlar har xil işorali: $a^2-4<0$, $a\in(-2;2)$.\\
— Karrali musbat ildiz: $D=4a+17=0$, $a=-\frac{17}4$, $y=\frac{15}4>0$ — mos.\\
$a\in\left\{-\frac{17}4\right\}\cup(-2;2)$, eng kiçigi $-\frac{17}4$.
\Javob{a) 2; \quad b) $-\frac{17}4$}

\newpage
\Savol{37}
$\cos x+\cos3x=2\cos2x\cos x$, \ $\sin x+\sin3x=2\sin2x\cos x$:
\[\cos2x\,(2\cos x+1)=\sin2x\,(2\cos x+1).\]
1) $\cos x=-\frac12$: $x=\pm\frac{2\pi}3+2\pi k$.\\
2) $\tg2x=1$ ($\cos2x=0$ da $\sin2x=\pm1$ — yeçim emas): $x=\frac\pi8+\frac{\pi k}2$.

a) Eng kiçik musbat: $\frac\pi8$.

b) $[-\pi;\pi]$: $\frac\pi8+\frac{\pi k}2$ dan $-\frac{7\pi}8,\,-\frac{3\pi}8,\,\frac\pi8,\,\frac{5\pi}8$ — $4$ ta; \ $\pm\frac{2\pi}3$ — $2$ ta. Jami $6$.
\Javob{a) $\frac\pi8$; \quad b) 6}

\newpage
\Savol{38}
$g$ grafigi $(0;3)$ dan ötadi $\Rightarrow$ $f(3)=0$. \ $f$ ning eng kiçik qiymati $f(-a)=b=-1$.\\
$(3+a)^2-1=0$ $\Rightarrow$ $a=-2$ yoki $a=-4$. \ $3\in[-a;\infty)$ bölişi kerak: $a=-4$ da $3<4$ — mos emas. Demak $a=-2$, $b=-1$:
\[f(x)=(x-2)^2-1,\ x\ge2;\qquad g(x)=2+\sqrt{x+1},\ x\ge-1.\]

a) $a+b=-3$.

b) $f$ ösuvçi, $g$ — uning teskarisi, kesişiş nuqtalari $y=x$ da yotadi:
\[(x-2)^2-1=x\ \Rightarrow\ x^2-5x+3=0\ \Rightarrow\ x=\frac{5\pm\sqrt{13}}2.\]
$\frac{5-\sqrt{13}}2\approx0{,}7<2$ — sohaga kirmaydi. Faqat $x=\frac{5+\sqrt{13}}2$: $1$ ta nuqta.
\Javob{a) $-3$; \quad b) 1}

\newpage
\Savol{39}
$f(x)=a(x-2)^2-3$, \ $f(0)=4a-3=-2$ $\Rightarrow$ $a=\frac14$:
\[f(x)=\frac14x^2-x-2,\quad a=\frac14,\ b=-1,\ c=-2.\]

a) $\dfrac{ac}b=\dfrac{-\frac12}{-1}=\dfrac12$.

b) $u=x-2$ deb olsak, parabola $y=\frac{u^2}4-3$, $A$ nuqta $u=2$, $y=2$. Masofa kvadrati:
\[d^2(u)=(u-2)^2+\left(\frac{u^2}4-5\right)^2,\quad \bigl(d^2\bigr)'=2(u-2)+u\left(\frac{u^2}4-5\right)=\frac{u^3-12u-16}4=\frac{(u-4)(u+2)^2}4.\]
Minimum $u=4$ da: nuqta $(6;1)$, \ $d=\sqrt{(6-4)^2+(1-2)^2}=\sqrt5$.
\Javob{a) $\frac12$; \quad b) $\sqrt5$}

\newpage
\Savol{40}
Siniq çiziq $(1;0)$ va $(5;0)$ dan ötadi, uçi $x=3$ da: $m=-3$, \ $-|1-3|+n=0$ $\Rightarrow$ $n=2$.\\
Parabola: $y=a(x-1)(x-5)$, uçi $(3;4)$: $a\cdot2\cdot(-2)=4$ $\Rightarrow$ $a=-1$:
\[y=-x^2+6x-5,\quad a=-1,\ b=6,\ c=-5.\]

a) $\dfrac{a+b-c}{n-m}=\dfrac{-1+6+5}{2+3}=2$.

b) Parabola ostidagi yuz minus uçburçak (asosi $4$, balandligi $2$):
\[S=\int_1^5\left(-x^2+6x-5\right)dx-\frac12\cdot4\cdot2=\frac{32}3-4=\frac{20}3.\]
\Javob{a) 2; \quad b) $\frac{20}3$}

\newpage
\Savol{41}
\fig{\begin{tikzpicture}[scale=0.62,font=\small]
\coordinate (A) at (0,0); \coordinate (C) at (8,0); \coordinate (B) at (5.5,4.330); \coordinate (D) at (5,0); \coordinate (I) at (5,1.732);
\draw (A)--(B)--(C)--cycle; \draw (I) circle (1.732);
\draw[axs] (I)--(D); \draw[ax] (I)--(C); \fill (I) circle (2pt);
\node[below left] at (A) {$A$}; \node[above] at (B) {$B$}; \node[below right] at (C) {$C$}; \node[below] at (D) {$D$}; \node[above] at (I) {$I$};
\node[aux,left] at (5,0.85) {$r$}; \node[below] at (6.5,0) {$3$};
\end{tikzpicture}}
$I$ — markaz. $CI$ — bissektrisa, $\angle ICD=30^\circ$, $ID\perp AC$.

a) $r=CD\cdot\tg30^\circ=\dfrac3{\sqrt3}=\sqrt3$.

b) Urinmalar xossasi: $CD=p-AB$ $\Rightarrow$ $p=3+7=10$.
\[S=p\cdot r=10\sqrt3.\]
\Javob{a) $\sqrt3$; \quad b) $10\sqrt3$}

\newpage
\Savol{42}
$AF\parallel BE$ — $ABEF$ trapetsiya, diagonallari $M$ da kesişadi. Trapetsiyada yon tomonlarga yopişgan uçburçaklar teng yuzli va
\[S_{ABM}=S_{EMF}=\sqrt{S_{AMF}\cdot S_{BME}}=\sqrt{32\cdot8}=16.\]
Xuddi şunday $EC\parallel FD$ — $ECDF$ trapetsiya, diagonallari $N$ da kesişadi:
\[S_{CND}=S_{ENF}=\sqrt{27\cdot3}=9.\]

a) $S_{MENF}=S_{EMF}+S_{ENF}=16+9=25$.

b) $S_{ABEF}=32+8+16+16=72$, \ $S_{ECDF}=27+3+9+9=48$: \ $S_{ABCD}=120$.
\Javob{a) 25; \quad b) 120}

\newpage
\Savol{43}
Oltiburçak tomoni $s$, \ $T=\frac{\sqrt3}4s^2$ — ikki tomoni $s$, orasidagi burçagi $120^\circ$ uçburçak yuzi (oltiburçak yuzining $\frac16$ qismi), masalan $S_{B_6B_5B_4}=T$.\\
$\angle B_1B_6B_5=120^\circ$, $\angle B_5B_6B_4=30^\circ$ $\Rightarrow$ $\angle B_1B_6B_4=90^\circ$, $B_6B_4=s\sqrt3$.\\
$\angle AB_6B_1=\varphi$ bölsa, $\angle EB_6B_4=90^\circ-\varphi$:
\[S_{AB_6B_1}=\frac12s^2\sin\varphi\cos\varphi=7,\qquad S_{EB_6B_4}=\frac12\left(s\sqrt3\right)^2\sin\varphi\cos\varphi=3\cdot7=21.\]
$S_{EB_6B_5B_4}=S_{EB_6B_4}-S_{B_6B_5B_4}$ $\Rightarrow$ $8=21-T$ $\Rightarrow$ $T=13$.

a) $B_6B_3=2s$ — katta diagonal, $B_2$ dan unga masofa $\frac{\sqrt3}2s$:
\[S_{B_6B_2B_3}=\frac12\cdot2s\cdot\frac{\sqrt3}2s=2T=26.\]

b) Oltiburçak va uni örab turgan töğri törtburçak umumiy simmetriya markaziga ega, şuning uçun $C$ va $D$ burçaklardagi qismlar mos ravişda $8$ va $7$ ga teng:
\[S_{ACDE}=6T+2\cdot7+2\cdot8=78+30=108.\]
\Javob{a) 26; \quad b) 108}

\newpage
\Savol{44}
\fig{\begin{tikzpicture}[scale=0.42,font=\small]
\coordinate (A) at (-6,0); \coordinate (D) at (6,0); \coordinate (C) at (4.82,1.82); \coordinate (B) at (-1.18,1.82); \coordinate (S) at (0,10.39); \coordinate (O) at (0,0); \coordinate (K) at (1.82,1.82);
\draw[dashed] (A)--(B)--(C); \draw[dashed] (S)--(B);
\draw (A)--(D)--(C); \draw (S)--(A); \draw (S)--(D); \draw (S)--(C);
\draw[ax] (S)--(O); \draw[ax] (O)--(B); \draw[ax] (O)--(C);
\fill (O) circle (3pt);
\node[below left] at (A) {$A$}; \node[below right] at (D) {$D$}; \node[right] at (C) {$C$}; \node[above left] at (B) {$B$}; \node[above] at (S) {$S$}; \node[below] at (O) {$O$};
\node[below] at (-3,0) {$6$}; \node[below] at (3,0) {$6$};
\end{tikzpicture}}
Asos $ABCD$: $AD=12$, $AB=BC=CD=x$, $\angle B=\angle C=120^\circ$ $\Rightarrow$ $\angle A=\angle D=60^\circ$, \ $AD=x+2\cdot\frac x2=2x$ $\Rightarrow$ $x=6$.\\
$AD$ ning örtasi $O$: $OA=OB=OC=OD=6$ ($\triangle OAB$, $\triangle OBC$, $\triangle OCD$ — muntazam). Yon qirralar asosga teng oğgan $\Rightarrow$ balandlik asosi — taşqi çizilgan aylana markazi $O$, $R=6$.\\
$SO=R\tg60^\circ=6\sqrt3$, \ yon qirra $SA=\frac R{\cos60^\circ}=12$.

a) $S_{\text{asos}}=3\cdot\frac{\sqrt3}4\cdot36=27\sqrt3$.\\
$SAD$: tomonlari $12$, $12$, $12$ — muntazam: $\frac{\sqrt3}4\cdot144=36\sqrt3$.\\
$SAB$, $SBC$, $SCD$: asosi $6$, yon tomonlari $12$, apofemasi $\sqrt{144-9}=3\sqrt{15}$: har biri $\frac12\cdot6\cdot3\sqrt{15}=9\sqrt{15}$.
\[S_{\text{töla}}=27\sqrt3+36\sqrt3+27\sqrt{15}=63\sqrt3+27\sqrt{15}.\]

b) $V=\frac13\cdot27\sqrt3\cdot6\sqrt3=162$.
\Javob{a) $63\sqrt3+27\sqrt{15}$; \quad b) 162}

\newpage
\Savol{45}
\fig{\begin{tikzpicture}[scale=0.9,font=\small]
\draw (0,0) circle (2); \coordinate (P) at (90:2); \coordinate (Q) at (180:2); \coordinate (R) at (-60:2);
\draw (P)--(Q)--(R)--cycle; \draw[axs] (0,0)--(P); \draw[axs] (0,0)--(Q); \draw[axs] (0,0)--(R); \fill (0,0) circle (1.5pt);
\node[aux] at (135:0.55) {\scriptsize$90^\circ$}; \node[aux] at (240:0.6) {\scriptsize$120^\circ$}; \node[aux] at (15:0.65) {\scriptsize$150^\circ$};
\node[above] at (P) {$P$}; \node[left] at (Q) {$Q$}; \node[below right] at (R) {$R$};
\node at (135:2.35) {$3$}; \node at (240:2.35) {$4$}; \node at (15:2.35) {$5$};
\end{tikzpicture}}
Berilgan perimetrda eng katta yuzali şakl — doira.

a) $2\pi r=12$ $\Rightarrow$ $r=2$ ($\pi\approx3$); \ $S=\pi r^2=3\cdot4=12$.

b) Yoylar $3$, $4$, $5$ — aylana uzunligining $\frac3{12}$, $\frac4{12}$, $\frac5{12}$ qismi, markaziy burçaklar $90^\circ$, $120^\circ$, $150^\circ$. Uçburçak markazdan çiqqan radiuslar bilan uçta uçburçakka bölinadi:
\[S=\frac12r^2\left(\sin90^\circ+\sin120^\circ+\sin150^\circ\right)=2\left(1+\frac{\sqrt3}2+\frac12\right)=3+\sqrt3.\]
\Javob{a) 12; \quad b) $3+\sqrt3$}

\ShowJavoblar
\end{document}
